\section{Pathogenic release strategies.} 

With both viral and bacterial particles being very small, it has historically been extremely challenging to make good observations of their behaviour. Progress is taking place, and in the last couple of decades, we have been able to observe individual bacteria through the course of an intracellular infection. For example, Dstl (Defence science technology laboratory) has been able to record videos of a macrophage bursting to release on the order of 100 Francisella Tularensis bacteria. Through the induction of bacterial luminescence, each individual pathogenic particle can be visually isolated and counted.\par
Francisella Tularensis and Y. Pestis are similar in size, each around 0.5 micrometers (Y. pestis being marginally larger). But viruses are generally about an order of magnitude smaller than bacteria, with HIV being approximately 100 nanometres. They are hence somewhat more elusive to visualisation and counting.\par 


\subsubsection*{Bursting}
Both Francisella Tularensis and Yersinia Pestis multiply within macrophages becfore causing the cell to lyse, releasing their pathogenic cohort for further replicative endevours. The case of Y. Pestis is slightly more complicated with a biphsaic switching; a change from intracellular to extracellular replication (more on that later). But generally, these strategies which involve a complete release in combination with the death of the cell are reffered to as bursting pathogenesis.

\subsubsection*{Budding}

As described above, budding release refers to the situation where an infected cell outputs new virions at a constant rate. Models of viral replication generally take this as an assumption and it has become a commonly held belief that this is how viral release takes place. 


\subsubsection*{Something in between?}

The recent work experimenting HIV release reveals that the story may not be so simple as was previously assumed. That in cases of non-lytic virulence there may be various strategies we are not fully aware of. Another example of something in between bursting and budding is release by extracellular vesicles (EVs). Amongst viruses that leave the cell non-lytically, there is a subdivision to be made into those which produce and assemble their own protein shell to harbour genetic material, and those which exploit the cell membrane to produce a lipid container. Of this second group, it has been known for a while that in some cases, multiple individual viruses get contained and released in one ex-membrane capsule. These are reffered to as extracellular vessicles \cite{kerviel2021new}, \cite{buzas2023roles}. They present another potential coplexity to the story of viral release.
